\section{结论}
\label{sec:conclusion}

本文提出并验证了一个面向基因编辑影响因素发现的科学发现平台，并以 DeepCRISPR 公开数据作为案例研究。平台的定位不是构造单一最优预测器，也不是证明表观遗传环境必然提升预测，而是把预测、模型归因、统计推断、环境增量分析、序列模式发现、细胞背景比较与证据整合组织为统一且可追溯的流程。

案例研究的主要结论是：（i）五类模型均能学到但仅能较弱地预测编辑效率（中位 $R^{2}$ 0.035--0.115），预测性能因此只作为进入解释阶段的准入条件；（ii）四个模型类独立地把 PAM 邻近种子区（17--20）识别为最高归因区域，第 18 位在跨模型平均谱中排序最高，其 C/A 效率差在 HCT116 与 HeLa 中约为 $+0.09$、在 HL60 中减小、在 HEK293T 中反向，构成一个具有上下文依赖性的候选序列关联；（iii）更宽的 CNN 感受野稳定改善预测（$k=7-k=3$：$+0.056$）并使归因向 PAM 邻近区域集中，而候选模式长度不随核大小改变，说明感受野不等于模式长度；（iv）环境因素提供的增量预测信息很小且不稳健（$|\Delta R^{2}|<0.01$，多数边级置信区间跨 0，区组析因方差分析未检出显著主效应）；（v）平台据此给出两个可执行的最小扰动候选（第 18 位碱基替换、PAM 邻近候选模式核心破坏），供后续实验验证。

综上，平台的科学性不来自"所有分析都得到阳性结果"，而来自能够在统一、可追溯的框架中区分稳定证据、模型特异性现象、上下文依赖现象与不确定结果，并把它们转换为可检验的生物学假设。后续工作将围绕两条主线展开：一是在独立编辑数据集上检验流程的可迁移性；二是对本文提出的候选假设开展最小扰动实验验证。
